% TOP_PROBLEM: {"id":30007616,"problem_number":"LOCAL-30007616","title":"Reserve-currency persistence","category_id":21,"category":{"id":21,"name":"economics","display_name":"Economics","description":"Open economic theory statements, published research agendas, and proposed scoped specifications; question provenance and historical priority are qualified per record.","slug":"economics","order_index":21,"created_at":"2026-10-08T00:00:00Z"},"status":"research_question","external_url":null,"legacy_ids":[],"published":false,"statement":"In the explicitly specified setting: Under what conditions would geopolitics, fiscal risk or digital money meaningfully weaken the dollar’s reserve-currency role? The mathematical statement above fixes the target, assumptions and scope of this catalogue formulation.","economics":{"original_id":105,"field":"International finance and exchange rates","kind":"theoretical","formalization_basis":"adapted_research_question","source_status":"research_agenda","priority_status":"earliest_found","priority_note":"This is the earliest explicit statement verified in this audit; first priority for the full catalogue formulation is not established.","earliest_verified_reference":{"authors":["Bank of England"],"title":"Bank of England Agenda for Research, 2025–2028","year":2026,"url":"https://www.bankofengland.co.uk/research/bank-of-england-agenda-for-research","doi":null,"locator":"Interconnected World / International finance, paragraph 2","evidence_summary":"Explicitly leaves evolution of global currencies open."},"current_reference":{"authors":["Bank of England"],"title":"Bank of England Agenda for Research, 2025–2028","year":2026,"url":"https://www.bankofengland.co.uk/research/bank-of-england-agenda-for-research","doi":null,"locator":"Interconnected World / International finance, paragraph 2","evidence_summary":"Explicitly leaves evolution of global currencies open."},"formalization_note":"Adapts an explicitly verified research-agenda passage. Mathematical objects, interventions, objectives and scope are proposed here; existing model-specific solutions are not relabeled unsolved."},"statusQualification":"Published question or research agenda verified at the cited locator. The mathematical specification is adapted for this catalogue and includes additional assumptions; source publication does not establish first-ever priority or certify this exact model remains unsolved. Adapts an explicitly verified research-agenda passage. Mathematical objects, interventions, objectives and scope are proposed here; existing model-specific solutions are not relabeled unsolved.","statusReviewedAt":"2026-10-08"}
% FIRST_POSTED: 2026-10-08
% ENTRY_KIND: statement_only
% !TeX program = lualatex
\documentclass[11pt]{article}
\usepackage{fontspec}
\usepackage[a4paper,margin=25mm]{geometry}
\usepackage{amsmath,amssymb,mathtools}
\usepackage{xurl}
\usepackage[hidelinks]{hyperref}
\newcommand{\sref}[2]{\hyperlink{source:#1:#2}{[S#2]}}
\setlength{\emergencystretch}{3em}
\begin{document}
% problemId:30007616
\section*{Reserve-currency persistence}\label{30007616}
\subsection{Definitions and mathematical statement}
Consider central banks choosing reserve portfolios across currencies \(j=1,\ldots,J\). Currency \(j\) offers returns, liquidity, settlement compatibility and exposure to issuer fiscal or sanctions risk. Let \(v_j\) be its equilibrium reserve share, with \(v_j\geq0\) and \(\sum_jv_j=1\). A reserve holder's payoff from currency \(j\) depends on fundamentals \(x_j\) and an explicitly specified network-benefit function \(n_j(v_j)\). Digital payment technologies alter settlement costs and feasible choices; these changes are primitives, not assumed conclusions. The task is to characterize equilibria and transition thresholds at which a persistent shock to \(x\) or payment technology substantially lowers the dollar share. Fix a material-decline threshold \(\eta\in(0,1)\) and horizon \(H\), and study
\[\Pr(v_{\mathrm{USD},t+H}\leq v_{\mathrm{USD},t}-\eta\mid do(x_t=x')).\]
Here USD denotes the dollar and \(x'\) is a fully specified intervention. Identify the role of coordination, switching costs and available substitutes, including multiple equilibria and hysteresis. Discipline network-benefit and risk parameters using reserve-holder behavior rather than assuming dominance is permanent. Distinguish reserve holdings from trade invoicing and private funding roles. Predicting one function's decline does not establish collapse of every international role of the currency.

\par\textbf{Formulation and scope.} Adapts an explicitly verified research-agenda passage. Mathematical objects, interventions, objectives and scope are proposed here; existing model-specific solutions are not relabeled unsolved.

\par\textbf{Open-question provenance.} An explicit question or research agenda was verified at the source locator. This does not by itself certify that every later version remains unresolved \sref{30007616}{1}.

\par\textbf{Historical priority.} This is the earliest explicit statement verified in this audit; first priority for the full catalogue formulation is not established.

\subsection{Short English statement}
In the explicitly specified setting: Under what conditions would geopolitics, fiscal risk or digital money meaningfully weaken the dollar’s reserve-currency role? The mathematical statement above fixes the target, assumptions and scope of this catalogue formulation.

\subsection{Sources}
\begin{itemize}\raggedright
\item[\textnormal{[S1]}]\hypertarget{source:30007616:1}{} Bank of England. 2026. Bank of England Agenda for Research, 2025–2028. Interconnected World / International finance, paragraph 2.
\url{https://www.bankofengland.co.uk/research/bank-of-england-agenda-for-research}.
{\small\textit{Earliest verified question/agenda reference; historical priority qualified below. Explicitly leaves evolution of global currencies open.}}
\end{itemize}
\end{document}
